\documentclass{article}
\usepackage[margin=1in]{geometry}
\usepackage{graphicx}
\usepackage{float}
\usepackage{microtype}
\usepackage{textcomp, gensymb}
\usepackage{enumitem}
\usepackage{amsmath}
\usepackage{parskip}
\begin{document}
\setcounter{section}{5}
\section{Chapter 6: Maxwell's Equations for Time-Varying Fields}
\setcounter{subsection}{-1}
\subsection{Dynamic Fields}
\(\vec{E}\) and \(\vec{D}\) in the static electric model \textbf{are not} related to \(\vec{H}\) and \(\vec{B}\) in the static magnetic model. In addition, static electric and magnetic fields \textbf{do not} give rise to waves that propagate and carry energy and information.

In this chapter, we will see that a \textbf{changing} magnetic field induces an electric field, and vice versa; that is, they are coupled, and \(\vec{E}\) and \(\vec{D}\) will be related to \(\vec{H}\) and \(\vec{B}\), contrary to their lack of interaction in the static model.

\subsection{Faraday's Law}

To review, in the static case, 
\[
  \nabla \times \vec{E} = 0.
\]

Modifying this equation to cover time-varying fields, in a general sense, 
\[
  \nabla \times \vec{E} = \frac{-\partial \vec{B}}{\partial t}.
\]
When \(\frac{\partial \vec{B}}{\partial t} = 0\), this relationship simplifies to the static case. Taking the surface integral of both sides of the equation,
\[
  \int \int_S (\nabla \times \vec{E}) \cdot d \vec{s} = - \int \int_S \frac{\partial \vec{B}}{\partial t} \cdot d \vec{s}.
\]
Using Stoke's theorem to simplify, 
\[
  \oint_C \vec{E} \cdot d \vec{l} = \frac{-d}{dt} \int \int \vec{B} \cdot d \vec{s} = \frac{-d}{dt} \ \Phi.
\]

The process described by the above equation is known as \textbf{electromagnetic induction}, and the left-hand side can be called the induced voltage, \(V_{\mathrm{emf}}\).

Extending to a general case, given \(N\) turns in a circuit, 
\[
	V_{\mathrm{emf}} = -N \ \frac{d\Phi}{dt}.
\]
In other words, the voltage induced by a changing magnetic field in a closed circuit is proportional to the rate of change in the magnetic flux and the number of turns \(N\) in the circuit. Obviously, the induced voltage then induces a corresponding current.

There are different types of induction. 1) stationary circuit/time-varying \(\vec{B}\) \(\to\) transformer induction. This is how transformers operate, by using a changing magnetic field to induce a voltage/current in a stationary circuit. 2) Moving circuit/static \(\vec{B}\) \(\to\) motional induction. 3) moving circuit/time-changing \(\vec{B}\) \(\to\) general case.

Mathetmatically, 
\[
	V_{\mathrm{emf}} = V_{\mathrm{emf}}^{\mathrm{tf}} + V_{\mathrm{emf}}^{\mathrm{m}}.
\]

\subsection{Stationary Loop in a Time-Varying Field}

Reviewing each in more detail, starting with transformer induction, recall that in this type of induction, there will be an induced current generated in a closed loop/circuit due to a changing \(\vec{B}\) field. Note that the flux produced by the induced current opposes the change in \(\vec{B}\) (Lenz's Law). 

Imagine a loop of wire (stationary) through which a changing \(\vec{B}\) is passing, and it is growing weaker with time. What direction will the induced current flow? Knowing that the \(\vec{B}\) field is decreasing, the induced \(\vec{B}\) field will \textit{oppose} this change, and increase. Using the right-hand rule, we can find the direction of the induced current. In the opposite case, where \(\vec{B}\) increases, the current will flow in the \textit{opposite} direction.

Example time. A circular loop of \(N\) turns of conducting wire lies in the \(x\)-\(y\) plane with its center at the origin of a magnetic field given by \(\vec{B} = \hat{\alpha_z} \ B_0 \cos{(\frac{\pi r}{2b})}\sin{(\omega t)}\), where \(b\) is the radius of the loop and \(\omega\) is the angular frequency. Given this, find the emf induced in the loop.

First, we can identify the type of induction as transformer induction, since the loop is stationary, and \(\vec{B}\) is a function of time \(t\). To find emf,
\[
	V_{\mathrm{emf}} = -N \frac{d\Phi}{dt}. 
\]
Therefore, we must first find magnetic flux \(\\Phi\) as a function of our circuit's structure and the magnetic field \(\vec{B}\).

Using the formula for magnetic flux,
\[
  \Phi = \int \int_S \vec{B} \cdot d \vec{s}.  
\]

In this case, for a single loop, \(d \vec{s} = \hat{\alpha}_z \ r \ dr \ d\phi\).

So, 
\[
	\Phi = \int_{0}^{2\pi} \int_{0}^{b} (\hat{\alpha}_z \ B_0 \cos{(\frac{\pi r}{2b})\sin{(\omega t)}}) \cdot (\hat{\alpha}_z \ r \ dr \ d\phi).
\]

Simplifying using integration by parts,
\[
	\Phi = 2\pi B_0 \sin{(\omega t)} \int_{0}^{b} \cos{(\frac{\pi r}{2b})} \cdot r \cdot dr = 2\pi B_0 \sin{(\omega t)}(\frac{2b^2}{\pi} - \frac{4b^2}{\pi^2}).
\]

Taking its time derivative, only the \(\sin\) function matters, so
\[
	\frac{d\Phi}{dt} = 2 \pi B_0(\frac{2b^2}{\pi} - \frac{4b^2}{\pi^2}) \omega \cos{(\omega t)}
\]

and so
\[
	V_{\mathrm{emf}} = -2 N \pi B_0(\frac{2b^2}{\pi} - \frac{4b^2}{\pi^2}) \omega \cos{(\omega t)} \ \mathrm{V}. 
\]

Another example problem: an inductor is formed by winding \(N\) turns of a thin conducting wire into a circular loop of radius \(a\). The inductor loop is in the \(x\)-\(y\) plane with its center at the origin, and connected to a resistor \(R\). In the presence of a magnetic field \(\vec{B} = B_0 (2 \hat{\alpha}_y + 3 \hat{\alpha}_z) \sin{(\omega t)}\). Find
\begin{enumerate}[label=\alph*)]
	\item the magnetic flux linking a single turn of the inductor.
		In this case, the changing magnetic field has both \(y\) and \(z\) components, and not all of it will end up as magnetic flux \(\Phi\) as a result. \(d \vec{s}\) is the same as in the prior problem, \(d \vec{s} = \hat{\alpha}_z \ (r \ dr \ d\phi)\). Finding magnetic flux using the same integral setup as before, 
		\[
			\Phi = \int_{0}^{2\pi}\int_{0}^{a} B_0 (2 \hat{\alpha}_y + 3 \hat{\alpha}_z)\sin{(\omega t)} \cdot (\hat{\alpha}_z \ (r \ dr \ d\phi)).
		\]
		Simplifying, 
		\[
			\Phi = 3B_0 \pi a^2 \sin{(\omega t)}.
		\]
	\item the emf given that \(N=10\), \(B_0 = a2T\), \(a=10 \ \mathrm{cm}\), and \(\omega=10^3 \ \mathrm{\frac{rad}{s}}\).

		Knowing that 
		\[
			V_{\mathrm{emf}} = -N \frac{d\phi}{dt} = -N \cdot 3 \pi a^2 B_0 \omega \cos{(\omega t)},
		\]
		substituting values, 
		\[
			V_{\mathrm{emf}} = -188.5 \cos{(10^3 t)} \ \mathrm{V}.
		\]
	\item the direction of the induced current at \(t=0\).
		Recall that induced current always opposes the direction of change in the magnetic flux; therefore, the question is: is magnetic flux \(Phi\) increasing or decreasing at \(t=0\)?

		We may evaluate the derivative \(\dfrac{d\Phi}{dt}\) at \(t=0\) to determine whether \(\Phi\) is increasing or decreasing. At \(t=0\), the derivative is positive (greater than zero), and so we may assume that \(\Phi\) is increasing at \(t=0\). Therefore, flux is increasing in the \(z\) direction at \(t=0\), and so the current will flow in the clockwise direction (right-hand rule for the negative \(z\) direction).
\end{enumerate}

\subsection{The Ideal Transformer}
An ideal transformer has two sides, with loops of wire (\(N\) turns) wrapped around each side; generally, one side is connected to an AC source, and the other to a load. Using Faraday's Law, 
\[
  V = -N \frac{d\Phi}{dt}, 
\]
we can generalize to the case of an ideal transformer, where
\[
	\frac{V_1}{V_2} = \frac{N_1}{N_2},
\]
while 
\[
  \frac{I_1}{I_2} = \frac{N_2}{N_1}.
\]
In addition,
\(R_{\mathrm{in}} = R_L \frac{N_1}{N_2}^2\) and \(Z_{\mathrm{in}} = Z_L \frac{N_1}{N_2}^2\).
  
\subsection{Moving Conductor in a Static Magnetic Field}
As an example, a conducting wire of length \(l\) moves across a static magnetic field \(\vec{B} = \hat{\alpha}_zB_0\) with velocity \(\vec{u}=\hat{\alpha}_x u\).

Recall that 
\[
  \vec{F}_m = q(\vec{u} \times \vec{B}).
\]
Therefore, 
\[
  \vec{E}_m = \frac{\vec{F}_m}{q} = \vec{u} \times \vec{B}.
\] 
This will induce a potential difference across the two ends of the wire, and 
\[
	V_{\mathrm{emf}} = \int_{2}^{1}\vec{E}_m \cdot d \vec{l} = \int_{2}^{1} (\vec{u} \times \vec{B}) \cdot d \vec{l} = \int_{}^{} (\hat{\alpha}_x u x \hat{\alpha}_z B_0) \cdot \hat{\alpha}_y dy = -u B_0 l.
\]

In general, if any segment of a closed circuit moves with a velocity \(\vec{u}\) across a static magnetic field \(\vec{B}\), then the induced motional emf is 
\[
	V_{\mathrm{emf}}^m = \oint (\vec{u} \times \vec{B}) \cdot d \vec{l}.
\]

As an example, a metal bar slides over a pair of conducting rails in a uniform magnetic field \(\vec{B} = \hat{\alpha}_z B_0\) with a velocity \(\vec{u} = \hat{\alpha}_x u\).
\begin{enumerate}[label=\alph*)]
	\item Determine the open circuit voltage \(V_0\) that appears across terminals 1 and 2.

		First, determine the direction of the current; current will flow from protons to electrons, so clockwise. Then, find the voltage,
		\begin{align*}
			V_0 &= \int_{}^{} (\vec{u} \times \vec{B}) \cdot d \vec{l} \\
					&=-u B_0 h \ \mathrm{V}
		\end{align*}
	\item Assuming that a resistance \(R\) is connected between the terminals, find the electric power dissipated in the resistor.

		First, 
		\[
		  I = \frac{u B_0 h}{R},
		\]
		and so
		\[
		  P = I^2 R.
		\]
\end{enumerate}
As another example, the rectangular loop shown below is located in the \(x\)-\(y\) plane and moves away from the origin with velocity \(\vec{u} = 5 \hat{\alpha}_y \mathrm{\frac{m}{s}}\) in a magnetic field given by \(\vec{B}(y) = \hat{\alpha}_z 0.2 e^{0.1y}\) \ mathrm{T}. If \(R = 5 \ohm\), find the current \(I\) at the instant that the loop sides are at \(y_1 = 2\) and \(y_2 = 2.5 \ \mathrm{m}\).

Recall the formula for motional emf, 
\[
	V_{\mathrm{emf}}^m = \oint_C (\vec{u} \times \vec{B}) \cdot d \vec{l}.
\]
Since \(\vec{u} \times \vec{B}\) is along \(\hat{\alpha}_x\), voltages are induced across only the sides oriented along \(\hat{\alpha}_x\) (1 and 2, and 3 and 4).
\end{document}
